\documentclass{beamer}
\usepackage{amsmath}
\usepackage{amssymb}

% Choose a theme
\usetheme{Madrid}
\usecolortheme{default}
\usefonttheme{serif}

% --- CUSTOM FOOTER ---
% Redefine the footline to remove author/institute and keep title/page number
\makeatletter
\setbeamertemplate{footline}
{
  \leavevmode%
  \hbox{%
  \begin{beamercolorbox}[wd=.5\paperwidth,ht=2.25ex,dp=1ex,center]{title in head/foot}%
    \usebeamerfont{title in head/foot}\insertshorttitle
  \end{beamercolorbox}%
  \begin{beamercolorbox}[wd=.5\paperwidth,ht=2.25ex,dp=1ex,right]{date in head/foot}%
    \usebeamerfont{date in head/foot}\insertframenumber{} / \inserttotalframenumber\hspace*{2ex}
  \end{beamercolorbox}}%
  \vskip0pt%
}
\makeatother
% --- END CUSTOM FOOTER ---

\title[Lecture 19: Waveform tomography]{Lecture 19: Waveform tomography}
\author{David Al-Attar}
\institute{Michaelmas term 2024}
\date{}

\begin{document}

% --- INTRODUCTION ---

\begin{frame}
    \titlepage
\end{frame}

\begin{frame}{Outline and Motivation}
    Today we discuss \textbf{waveform tomography}, a modern and powerful successor to the delay-time tomography method from the last lecture.
    \begin{itemize}
        \item Unlike delay-time methods which use only arrival times, waveform tomography uses the \textbf{entire seismogram}.
        \pause
        \item It fully accounts for the complex physics of wave propagation by using numerical solutions to the elastic wave equation.
        \pause
        \item It treats the inverse problem as fundamentally \textbf{non-linear}.
        \pause
        \item Our focus will be on how to solve this large-scale inverse problem using \textbf{gradient-based optimisation}.
    \end{itemize}
\end{frame}

% --- METHOD ---

\begin{frame}{Waveform Misfit}
    The first step is to quantify the difference between an observed seismogram, $u^{obs}(t)$, and a synthetic seismogram, $u(t)$, calculated in a given Earth model.
    \pause
    \begin{block}{Least-Squares Misfit Function}
        A simple choice is the squared L2-norm of the difference between the waveforms:
        $$ J = \frac{1}{2}\int_{0}^{T} ||\mathbf{u}(\mathbf{x}_{r},t)-\mathbf{u}^{obs}(t)||^{2}dt $$
    \end{block}
    \pause
    \begin{itemize}
        \item The value of the misfit $J$ depends on the synthetic wavefield $\mathbf{u}$, which in turn depends on the underlying Earth model parameters, $\mathbf{m}$. We can therefore think of the misfit as a function $J(\mathbf{m})$.
        \pause
        \item \textbf{The goal of waveform tomography is to find the model $\mathbf{m}$ that minimises this misfit function.}
    \end{itemize}
\end{frame}

\begin{frame}{Solving the Inverse Problem: Optimisation}
    We want to find a model $\overline{\mathbf{m}}$ that is a minimum of the misfit function $J(\mathbf{m})$. This occurs where the gradient (or functional derivative) of the misfit is zero:
    $$ DJ(\overline{\mathbf{m}}) = 0 $$
    \pause
    \begin{itemize}
        \item The model space is extremely high-dimensional, so methods like grid search or random sampling (MCMC) are computationally impossible.
        \pause
        \item We must use \textbf{local, gradient-based optimisation methods}. These iteratively improve an initial model by moving it in a direction related to the negative gradient.
        \pause
        \item A major drawback is that these methods can get stuck in a \textbf{local minimum} of the misfit function, which may not be the true global minimum.
    \end{itemize}
\end{frame}

\begin{frame}{The Steepest Descent Algorithm}
    The simplest gradient-based method is steepest descent. It is an iterative process:
    \pause
    \begin{enumerate}
        \item Start with an initial guess for the model, $\mathbf{m}_0$.
        \pause
        \item At the current model $\mathbf{m}_k$, calculate the gradient of the misfit, $DJ(\mathbf{m}_k)$. The direction of steepest descent is $-\,DJ(\mathbf{m}_k)$.
        \pause
        \item Update the model by taking a step in that direction:
        $$ \mathbf{m}_{k+1} = \mathbf{m}_k - \lambda \, DJ(\mathbf{m}_k) $$
        \pause
        \item Find the optimal step length $\lambda$ using a line search to achieve the largest misfit reduction.
        \pause
        \item Repeat until the gradient is nearly zero or the misfit stops decreasing.
    \end{enumerate}
\end{frame}

% --- ADJOINT METHOD ---

\begin{frame}{Calculating the Gradient}
    The key challenge is calculating the gradient $DJ(\mathbf{m})$ efficiently.
    \begin{block}{A Bad Idea: Finite Differences}
        We could approximate the derivative by perturbing each model parameter one by one and running a new simulation for each.
        $$ \frac{\partial J}{\partial m_i} \approx \frac{J(\mathbf{m} + \epsilon \mathbf{e}_i) - J(\mathbf{m})}{\epsilon} $$
        This would require $N+1$ wavefield simulations, where $N$ is the number of model parameters. Since $N$ can be millions, this is computationally prohibitive.
    \end{block}
    \pause
    \begin{alertblock}{A Good Idea: The Adjoint Method}
        There is a much more efficient way to calculate the exact gradient, known as the \textbf{adjoint method}.
    \end{alertblock}
\end{frame}

\begin{frame}{The Adjoint Method (1/2): Concept}
    The adjoint method uses the method of Lagrange multipliers to find the gradient of a function ($J$) subject to a constraint (the wave equation).
    \begin{itemize}
        \item The derivation is lengthy, but the conceptual result is elegant.
        \pause
        \item It introduces an \textbf{adjoint wavefield}, $u'$. This adjoint field satisfies the same wave equation as the normal "forward" wavefield.
        \pause
        \item However, the source for the adjoint wavefield is the \textbf{data residual} ($u - u^{obs}$), but applied as a force at the \textit{receiver} location and running \textbf{backwards in time}.
    \end{itemize}
\end{frame}

\begin{frame}{The Adjoint Method (2/2): The Payoff}
    \begin{block}{Calculating the Sensitivity Kernel}
        The gradient of the misfit with respect to any model parameter can be calculated by combining the forward field ($\mathbf{u}$) and the adjoint field ($\mathbf{u}'$). For example, the sensitivity to density, $\rho$, is:
        $$ K_{\rho} = \int_{0}^{T} \frac{\partial \mathbf{u}}{\partial t} \cdot \frac{\partial \mathbf{u}'}{\partial t} dt $$
        Similar expressions exist for elastic parameters.
    \end{block}
    \pause
    \begin{alertblock}{Computational Advantage}
        The total cost to compute the gradient for \textbf{all} model parameters is:
        \begin{enumerate}
            \item One forward simulation to get $\mathbf{u}$ and the data residual.
            \item One backward (adjoint) simulation to get $\mathbf{u}'$.
            \item One step to combine the saved fields.
        \end{enumerate}
        The cost is independent of the number of model parameters, making large-scale tomography feasible.
    \end{alertblock}
\end{frame}

% --- BANANA DOUGHNUT KERNELS ---

\begin{frame}{Banana-Doughnut Kernels}
    The sensitivity kernels calculated with the adjoint method have a particular structure.
    \begin{itemize}
        \item Let's consider the sensitivity of a travel-time measurement to wave speed (a case related to both delay-time and waveform tomography).
        \pause
        \item Ray theory states that the sensitivity is confined to an infinitely thin geometric ray path.
        \pause
        \item The full-wave theory shows that the sensitivity is actually spread out over a volumetric region around the ray path. This region is the first \textbf{Fresnel zone}.
        \pause
        \item Curiously, the sensitivity is zero exactly on the geometric ray itself.
    \end{itemize}
    \pause
    \begin{center}
    \textit{This shape -- broad around the ray path but hollow in the middle -- led to the name \textbf{"banana-doughnut kernel"}.}
    \end{center}
\end{frame}

% --- SUMMARY ---

\begin{frame}{Summary: What You Need to Know}
    \begin{itemize}
        \item \textbf{Waveform Tomography} aims to minimise the misfit between observed and synthetic seismograms, using the full waveform.
        \pause
        \item The inverse problem is highly non-linear and solved with iterative, \textbf{gradient-based optimisation} methods like steepest descent.
        \pause
        \item The gradient of the misfit function is calculated efficiently using the \textbf{adjoint method}, which involves one forward and one backward (adjoint) wave simulation.
        \pause
        \item The cost of the adjoint method is independent of the number of model parameters, making large-scale problems computationally tractable.
        \pause
        \item Full-wave sensitivity kernels (\textbf{banana-doughnut kernels}) have finite width related to the Fresnel zone, differing from the infinitely thin sensitivities of ray theory.
    \end{itemize}
\end{frame}

\end{document}
